\cleardoubleoddpage
\chapter{Introduction}
\label{chap:introduction}

\section{Motivation}

Magnetically confined fusion plasmas exhibit turbulence that transports heat and particles across
magnetic surfaces. Predicting this transport is important for understanding confinement, designing
experiments, and deciding which reduced models are safe to use. Nonlinear gyrokinetic equations
retain kinetic effects that fluid approximations can omit, but they evolve a distribution function
over several spatial and velocity coordinates. The resulting five-dimensional state is expensive to
advance at high fidelity, and a long simulation produces a large sequence of strongly correlated
snapshots \citep{abel2013multiscale,paischer2025gyroswin}.

The computational cost creates an attractive role for learned surrogates. A surrogate could provide
fast approximate diagnostics between expensive solver calls, help explore parameter settings, or
serve as a proposal model for a larger simulation workflow. These applications require more than a
low training loss. The model must preserve the information needed by a physical diagnostic, respect
the temporal ordering of trajectories, and be compared with a baseline that is meaningful for the
quantity of interest. If those conditions are not met, a visually smooth rollout can be less useful
than copying the last observed diagnostic.

This thesis studies a deliberately narrow version of that problem. Rather than reconstructing the
complete distribution function, it asks whether a snapshot can be compressed into a small latent
state whose temporal dynamics are easier to learn while still supporting scalar-flux and spectral
diagnostics. Let $x_t$ denote a preprocessed channel-first snapshot, $z_t=f_\theta(x_t)$ its
latent representation, and $g_\phi$ the diagnostic heads. A temporal model $h_\psi$ predicts future
states from a context window. The complete prediction route is therefore
\begin{equation}
  x_{t-L+1:t}\;\xrightarrow{f_\theta}\;z_{t-L+1:t}
  \xrightarrow{h_\psi}\widehat{z}_{t+1:t+H}
  \xrightarrow{g_\phi}\widehat{y}_{t+1:t+H},
\end{equation}
where $y$ contains flux and spectra. The route separates representation learning, temporal
prediction, and diagnostic decoding so that each source of error can be inspected.

This separation has practical advantages. Sequence models operate on vectors rather than a full
spatial grid, so their memory use is easier to control. Diagnostic heads make it possible to test
whether a latent state retains task-relevant information without pretending that it is a complete
physical field. A cached representation also allows persistence, GRU, and Transformer baselines to
share exactly the same encoder and target normalization. The trade-off is equally important: a
latent state is useful only for the information it retains. A small latent MSE does not imply an
accurate transport forecast, and an accurate scalar diagnostic does not imply full-field fidelity.

\section{Why Evaluation Design Matters}

Trajectory data make leakage easy to introduce. Adjacent snapshots from one trajectory are highly
correlated, and a random snapshot split can place near-duplicates in training and evaluation. More
subtly, an encoder trained with one trajectory split can be combined with a sequence evaluator
using another split. The downstream evaluator may call a trajectory ``test'' even though its
snapshots already influenced the encoder or latent normalization. The resulting number can be
finite and reproducible while failing to measure end-to-end generalization.

The thesis therefore treats provenance as part of the experiment, not as an appendix. A manifest
binds the ordered trajectory universe and each outer fold. Resolved configurations record data and
training seeds, normalization scope, context, horizon, and upstream artifact hashes. Validation
selects checkpoints and model families; outer-test evidence is opened only after the selection
barrier. A direct observed-flux persistence baseline is declared before comparison because the
scientific question is diagnostic prediction, not merely latent reconstruction. This design makes
a negative result interpretable: if a learned model loses to persistence, the loss is not hidden by
a weaker reference or by test-aware tuning.

\section{Research Question}

The central research question is:

\begin{quote}
Can a compact JAX/Flax latent surrogate learn short-horizon gyrokinetic dynamics from
preprocessed five-dimensional field snapshots, and do latent sequence models improve
flux-predictive rollouts over direct persistence of the observed flux on held-out trajectories?
\end{quote}

The phrase ``held-out'' is used at the outer-fold level. Because the 51-trajectory universe was
available during development, the final estimate is retrospective nested group cross-validation,
not a pristine locked test. Three sub-questions make the central question measurable:

\begin{enumerate}[label=RQ\arabic*]
  \item Does the encoder retain enough diagnostic signal in a 128-dimensional state for finite
        flux and spectral prediction?
  \item After validation-only checkpoint and family selection, does the selected learned sequence
        model beat direct persistence of the observed flux over an eight-step horizon?
  \item Is any apparent learned-model advantage stable across matched training seeds and outer
        folds, or does the selected family change with the split?
\end{enumerate}

Flux RMSE is the primary metric because transport prediction is the stated goal. Latent MSE,
flux MAE, spectra errors, shape correlation, and finite-output stability are secondary diagnostics.
The hierarchy is intentional: a model that predicts a latent vector well but decodes a poor flux is
not a successful answer to the central question.

\section{Contributions}

This thesis makes the following contributions:

\begin{itemize}
  \item a portable JAX/Flax implementation of a latent-surrogate workflow, including synthetic
        fixtures and a real Cyclone/KvikIO data adapter;
  \item an explicit channel-first tensor contract for five spatial dimensions and deterministic
        trajectory-level split manifests;
  \item a representation objective combining SimSiam-style consistency with supervised flux and
        spectra heads;
  \item latent caching and comparable latent-state persistence, GRU, and causal Transformer
        candidates;
  \item trajectory-balanced rollout evaluation with paired hierarchical uncertainty and separate
        spectral diagnostics;
  \item a frozen five-fold, five-seed evidence record with 230 accepted and 25 skipped stage slots;
        and
  \item a reproducibility-oriented CLI, configuration system, test suite, and release manifest
        that make the negative result auditable without publishing raw data or private paths.
\end{itemize}

The empirical contribution is deliberately bounded. The accepted result shows that the selected
latent models do not beat observed-flux persistence under the declared data and horizon. It does
not establish that latent surrogates are unsuitable for gyrokinetics, nor does it compare numerical
performance with GyroSwin. The software and protocol provide a controlled basis for that future
question.

\section{Scope and Claim Boundaries}

The work predicts latent vectors and diagnostics; it does not reconstruct the complete
five-dimensional state. The medium universe contains 51 trajectories from one preprocessing
pipeline, and the retrospective outer folds provide ten or eleven test trajectories each. The
eight-step horizon establishes short-term behaviour only. Generalization to other regimes,
machines, preprocessing choices, or substantially longer horizons is not established.

GyroSwin is discussed as related work because it targets scalable full-field five-dimensional
surrogacy \citep{paischer2025gyroswin}. A direct numeric comparison is intentionally omitted. Such
a comparison would require the same dataset revision, trajectory split, temporal sampling,
forecast horizon, target definition, normalization, checkpoint, and metric implementation. Without
that matched bundle, a table of numbers would be misleading rather than informative.

\section{Thesis Structure}

\autoref{chap:background} introduces gyrokinetic turbulence, representation learning, and latent
sequence models. \autoref{chap:method} defines the data contract, objectives, architectures, and
evaluation protocol. \autoref{chap:implementation} describes the software system and portability
decisions. \autoref{chap:experiments} specifies the byte-verified universe, nested folds, models,
metrics, and provenance rules. Results are reported in \autoref{chap:results}, interpreted in
\autoref{chap:discussion}, and summarized in \autoref{chap:conclusion}. The required story summary
and reproducibility details are provided in the appendices.
